%%%%%%%%%%%%%%%%%%%%%%%%%%%%%%%%%%%%%%%%%%%%%%%%%%%%%
%			GENERÁTOR 18.09  						%
%%%%%%%%%%%%%%%%%%%%%%%%%%%%%%%%%%%%%%%%%%%%%%%%%%%%%
\documentclass[openany,12pt]{memoir}
\usepackage[utf8]{inputenc} 
\usepackage[czech]{babel}
\usepackage[T1]{fontenc}
\usepackage[top=1.5cm, bottom=2cm, left=2cm, right=2cm]{geometry}  % --> NASTAVENÍ OKRAJŮ
\usepackage{fancyhdr}
\usepackage{graphicx}
\usepackage{xwatermark}
\usepackage{xcolor}
\usepackage{changepage}
\usepackage{pdfpages}
\usepackage{lettrine}
%\chapterstyle{veelo}  %neni potřeba

%%%%%%%%%%%%%%%%%%%%%%%%%%%%%%%%%%%%%%
%  FONT                              %
%%%%%%%%%%%%%%%%%%%%%%%%%%%%%%%%%%%%%%
\usepackage[charter]{mathdesign}



%%%%%% Package na zpěvník
\usepackage[full]{leadsheets}%http://mirrors.nic.cz/tex-archive/macros/latex/contrib/leadsheets/leadsheets_en.pdf   --> dokumentace	
\definesongtitletemplate{empty}{} 
\setchords{
format = \bfseries,   %tučné akordy
minor = {mi},% 
input-notation = {german},%
output-notation = {german}%
}
\definesongtitletemplate{empty}{} 

\newlength{\drop}
\newwatermark[pages=3-,color=red!50,angle=0,scale=2, xpos=0,ypos=0]{\includegraphics[width=5cm]{obr/pozadi2.jpg}} %--> dvojka na pozadí


%%%% Vlastní příkazy
\newcounter{Slokočet}   %Automatické číslování slok
\newcommand{\mezera}{\vspace*{0.5cm}}   %Horizontální odsazení slok
\newcommand{\stred}{5.2cm}   %%% Na zarovnání slok doprostřed, pozn. automatičtější zarovnávání na střed nejde
\newcommand{\refren}{\mezera \noindent \textbf{R:} } %refrén
\newcommand{\sloka}{\mezera \noindent \addtocounter{Slokočet}{1} \arabic{Slokočet}. } 	%sloka, která se automaticky čísluje
\newcommand{\ssloka}{\mezera \noindent}  % vlastní číslo sloky
\newcommand{\predehra}{\mezera \hspace*{-1.85cm} \textbf{Předehra:} }
\newcommand{\carka}{,\:}
\newcommand{\m}[1]{\color{white}{#1}}  %Pro akordy
\newcommand{\ap}{'}	%Pro apostrof
\newcommand{\elipsa}{\kern\fontdimen3\font} %Příkaz pro lepší zacházení s výpustkami (=...); je to vpodstatě jen mezera mezi tečkama výpustky

\addto\captionsczech{\renewcommand{\contentsname}{Seznam písní}}

%%%%%%%%%%%%%%%%%%%%%%%%%%%%%%%%%%%%
%    FORMÁTOVÁNÍ                   %
%%%%%%%%%%%%%%%%%%%%%%%%%%%%%%%%%%%%

%%% Vlevo zarovnaný text s blokem zarovnaným na střed
\usepackage{varwidth}% http://ctan.org/pkg/varwidth
\newenvironment{centerleft}{%
  \begin{center} % so the minipage is centered
  \begin{varwidth}[t]{\textwidth}	
  \raggedright % so the minipage's text is left justified
  \setlength{\parindent}{0.45cm}
}{%
  \end{varwidth}
  \end{center}
}

%%% Vícesloupcové stránky
\usepackage{multicol}



%%%%%%%%%%%%%%%%%%%%%%%%%%%%%%%%%%%%%%%%%%%%%%%%%%%%%%
%			NÁVOD									 %
%%%%%%%%%%%%%%%%%%%%%%%%%%%%%%%%%%%%%%%%%%%%%%%%%%%%%%
%1. Věci v hlavičce IGNOROVAT
%2. Píseň psát do prostoru mezi \begin{song} a \end{song}
%3. další řádek se značí dvěma odsazeníma (= dvakrát stisknout enter)
%4. \refren vždy na začátku refrenu a \sloka na začátek sloky (automaticky se čísluje)
%  \ = alt gr + q ; [/] = alt gr f/g ; {/} = alt gr + b/n; ^ = alt gr + 3 + mezera
%Cokoliv napíšete do ^{  } se bude brát jako akord
%když se toto bude dotýkat nějakého slova (nebude mezi tím a slovem mezera)
%tak se akordy zjeví nad slovem, ale pište to před slovo
%Když se to nedotýká slova, tak akord lítá ve vzduchu a vytiskne se větší mezera
%První možnost je asi preferovanější
%5. Akordy stačí psát jen do první sloky, když se nezmění -- kytaristi to zvládnou



%
%\usepackage{subfiles}
%</++>

\begin{document}
\sffamily %Sans-serif pro lepší čitelnost

\pagestyle{empty}
\newgeometry{top=1.5cm, bottom = 0cm, left = 2cm, right = 2cm}
	
\newpage

%Pro kompilaci je důležitý, aby ten \input a \end{document} zůstali na posledních
%dvou řádcích!!!!

\newpage
\begin{song}{title=\centering Název písně \\\normalsize Autor písně }  %% sem se napíše jméno songu a autor

%\vspace*{.5cm}

%\begin{centerjustified}
%\begin{varwidth}{\textwidth}
%madafaka
%\end{varwidth}
\begin{varwidth}{\textwidth}
\begin{multicols*}{2}
%\normalsize \vspace*{-2\baselineskip} %snotyzaklad
\normalsize \vspace*{-1\baselineskip} \vspace*{-1em} \vspace*{.02cm}%beznotzaklad

\sloka  %1
Sloka ^{E}nana na nanananana

^{Ami}Nové odsazení ^{D}lalala lalalallaa

\refren 

/: Refrén v ^{Ami}repetici

^{D}jak chytlavé\elipsa.\elipsa.\elipsa. :/

\sloka %2
Sloka ^{E}nana na nanananana

^{Ami}Nové odsazení ^{D}lalala lalalallaa

\refren 

/: Refrén v ^{Ami}repetici

^{D}jak chytlavé\elipsa.\elipsa.\elipsa. :/

\sloka %3
Sloka ^{E}nana na nanananana

^{Ami}Nové odsazení ^{D}lalala lalalallaa

\refren 

/: Refrén v ^{Ami}repetici

^{D}jak chytlavé\elipsa.\elipsa.\elipsa. :/

\sloka %4
Sloka ^{E}nana na nanananana

^{Ami}Nové odsazení ^{D}lalala lalalallaa

\refren 

/: Refrén v ^{Ami}repetici

^{D}jak chytlavé\elipsa.\elipsa.\elipsa. :/

\sloka %5
Sloka ^{E}nana na nanananana

^{E}Nové odsazení lalala lalalallaa

\refren 

/: Refrén v ^{Ami}repetici

^{D}jak chytlavé\elipsa.\elipsa.\elipsa. :/

\sloka 
Sloka ^{E}nana na nanananana

^{Ami}Nové odsazení ^{D}lalala lalalallaa

\refren 

/: Refrén v ^{Ami}repetici

^{D}jak chytlavé\elipsa.\elipsa.\elipsa. :/

\sloka 
Sloka ^{E}nana na nanananana

^{Ami}Nové odsazení ^{D}lalala lalalallaa

\refren 

/: Refrén v ^{Ami}repetici

^{D}jak chytlavé\elipsa.\elipsa.\elipsa. :/
\hfill
\end{multicols*}
\end{varwidth}
\begin{varwidth}{\textwidth}
yeah boi
\end{varwidth}

%\end{centerjustified}


\setcounter{Slokočet}{0}
\end{song}



\end{document}

